% Printed212; principalPDF219.
\par\clearpage
\noindent212\hfill\textsc{Correspondence.}\hfill\mbox{}\par\medskip
\begin{center}Concise Sample Calculation.\\Neuwerk.\end{center}
\[
\begin{array}{r@{\qquad}r@{\qquad}r}
-266575{,}038&+95076{,}254&\\[8pt]
\begin{array}[t]{r}
133{,}287\,519\\10{,}663\,002\\\hline
143{,}950\,521\\2^\circ\,23'\,57''{,}03\\51\quad23\quad20{,}61\\\hline
53\quad47\quad17{,}64\\8\quad15{,}73\\\hline
53\quad55\quad33{,}37\\-31{,}40\\\hline
53^\circ\,55'\quad1''{,}97\\\hline
\end{array}&
\begin{array}[t]{r}
4{,}978\,0721\\8{,}508\,8620\\\hline
3{,}486\,9341\\9{,}769\,9904\\0{,}137\,5588\\9{,}956\,1442\\1{,}403\,2842\\\hline
3{,}716\,9437\\-764\\\hline
1{,}496\,9872\\-850\\\hline
3{,}624\,4929\\-923
\end{array}&
\begin{array}[t]{r}
7{,}43389\\\hline5{,}44934\\5{,}49566\\5{,}53146\\[8pt]
5210''{,}355\\1^\circ\,26'\,50''{,}355\\\hline
4211''{,}149\\1^\circ\,10'\,11''{,}15\\\hline
\end{array}
\end{array}
\]
\bigskip
\begin{center}\textsc{Gauss} to \textsc{Schumacher}.\quad Göttingen, 30 April 1830.\end{center}
. . . . . . I now continue with the auxiliary tables. Today I give only the formulas for the remaining three columns used in calculating longitude, latitude and convergence. I do not know, however, whether I denoted them by \(D,E\) and \(F\) in my previous letter.

Numbers whose logarithms are entered:
\[
\begin{array}{@{}l@{\;}l@{}}
\text{Longitude . . . . }[D=]&h\left(\frac34-\frac14\cos2\varphi+\frac{ee}{2(1-ee)}\cos\varphi^4\right)\\[6pt]
\text{Latitude . . . . }[E=]&h\left(\frac34+\frac{2-11ee}{4(1-ee)}\cos\varphi^2+\frac{5ee}{2(1-ee)}\cos\varphi^4-\frac{e^4}{(1-ee)^2}\cos\varphi^6\right)\\[6pt]
\text{Convergence }[F=]&h\left(1-\frac{ee}{1-ee}\cos\varphi^4-\frac{2e^4}{(1-ee)^2}\cos\varphi^6\right).
\end{array}
\]
\[
\text{Here we have }h=\frac{k\,.\,10^7}{3\,.\,(206265)^2};\qquad
\begin{array}{l}
k=\text{modulus}=0{,}434\rlap{\raisebox{-.5ex}{\tiny.}}^{\prime}2945\\
\log h=5{,}531\,8128-10.
\end{array}
\]

% Printed213; principalPDF220.
\par\clearpage
\noindent\mbox{}\hfill{\small\textsc{On the Formulas for the Hanoverian Land Survey.}}\hfill213\par\medskip
Note that these formulas are complete, that is, they are not infinite series and contain only these terms. . . . . . .

To increase your confidence in \textsc{Schmidt}'s calculation, I note that he calculated the two principal elements of the Earth's dimensions four times, but only once did he recalculate because of an error in arithmetic. Specifically:
\begin{list}{}{\setlength{\leftmargin}{2em}\setlength{\itemsep}{0pt}\setlength{\parsep}{0pt}}
\item[1)] The numbers in my determination of latitude, etc.

These had contained an error in arithmetic, which he later corrected, hence
\item[2)] the numbers in his Geography and in your A. N.

Only later did I draw his attention to the correction of the standards of length used in the East Indies, hence
\item[3)] the calculation whose result is in the preface to the book.

Finally, he has since made a fourth calculation. He did so to include among the data the results of \textsc{Struve}'s arc measurement, which had only subsequently come to his attention, and not on account of an error in arithmetic. The result
\item[4)] was communicated to me by him in manuscript and is the same as that underlying my new auxiliary tables, namely
\end{list}
\[
\text{Flattening}=\frac1{297{,}732};\qquad
\frac{\text{Earth's meridian quadrant}}{90}=57008{,}551.
\]
. . . . . .
\bigskip
\begin{center}\textsc{Gauss} to \textsc{Schumacher}.\quad Göttingen, 17 May 1831.\end{center}
. . . . . . There are \textls[120]{several} ways of calculating the coordinates from latitudes and longitudes, each with its own advantages in the appropriate situation. For the case in which you need such calculations, it is most convenient to use the series, which then converge so rapidly that very few terms suffice.

If the coordinates \(x,y\) correspond to latitude \(\varphi\) and longitude \(\lambda\), the form is as follows\textsuperscript{[*)]}:
\[
\begin{array}{l}
x=A-A'\lambda\lambda-A''\lambda^4-\text{etc.}\\
y=B\lambda+B'\lambda^3+\text{etc.},
\end{array}
\]
\par\noindent\rule{4cm}{.4pt}\par
{\footnotesize\noindent\hspace*{\parindent}[*) Cf. Art. 13, p. 158.]\par}

% Printed214; principalPDF221.
\par\clearpage
\noindent214\hfill\textsc{Correspondence.}\hfill\mbox{}\par\medskip
\noindent where \(A,A',A''\), etc., and \(B,B'\), etc., are functions of \(\varphi\). Their numerical values are to be calculated for the whole degrees (or half degrees, etc.) that occur within the map.

You need no instructions from me for calculating \(A\), since it depends solely on the rectification of an elliptic arc. Let \(\varphi^\circ\) be the latitude of the place from which the coordinates \(x\) are counted southward, \(u^\circ\) its actual distance from the equator, and \(u\) the variable distance of the parallel \(\varphi\) from the equator (both \(u,u^\circ\) expressed in the unit chosen for the coordinates). Then \(A=u^\circ-u\). For this you can also use an auxiliary table I sent you earlier\textsuperscript{[*)]}. You can calculate the values of the other coefficients using the following formulas, which are \textls[120]{absolutely} exact.
\[
\begin{aligned}
A'&=\frac{asc}{2p}\\[4pt]
A''&=\frac{asc}{24(1-ee)^2p}\bigl(5-ee-(6+6ee)ss+(9ee+3e^4)s^4-4e^4s^6\bigr)\\[4pt]
B&=\frac{ac}{p}\\[4pt]
B'&=\frac{ac}{6(1-ee)p}(1-2ss+ees^4).
\end{aligned}
\]
The following points should be noted.

The notation is as follows:
\begin{quote}
\(a\) denotes the radius of the equator and \(e\) the eccentricity;
\[
\GaussRightBrace{\begin{array}{l}
c=\cos\varphi\\s=\sin\varphi\\p=\surd(1-eess)
\end{array}}\ \text{for brevity;}
\]
\end{quote}
\noindent moreover, the unit taken for \(\lambda\) is the arc \(57^\circ17'45''\). If, for instance, you ultimately wish to calculate the coordinates for successive degrees of \(\lambda\), it is best at the outset to give the coefficient \(B\) the factor \(\frac1{57{,}296}=\frac{3600}{206265}=\frac\pi{180}\) and likewise to multiply the coefficients \(A',B',A''\) at once by the second, third and fourth powers of \(\frac\pi{180}\).

The formulas become somewhat simpler if you set \(\frac{ee}{1-ee}=m\) and recast them accordingly. I myself used them in this form for my
\par\noindent\rule{4cm}{.4pt}\par
{\footnotesize\noindent\hspace*{\parindent}[*) Page 206.]\par}

% Printed215; principalPDF222.
\par\clearpage
\noindent\mbox{}\hfill{\small\textsc{On the Formulas for the Hanoverian Land Survey.}}\hfill215\par\medskip
\noindent calculations, but did not keep them\textsuperscript{[*)]}. I therefore leave the very small task of that transformation to you. Taking \(\varphi^\circ=51^\circ31'47''{,}85\), using Mr. \textsc{Schmidt}'s latest determination of the flattening, and choosing the \(10\,000\,000\)th part of the Earth's meridian quadrant in this shape as the unit of length, I find, with \(\lambda=n\) degrees:
\[
\begin{array}{c@{\quad}l@{\qquad}l}
\varphi&x=&y=\\
51^\circ&+58947{,}1-475{,}95nn-0{,}0167n^4&70180{,}0n-0{,}737n^3\\
52^\circ&-52287{,}0-472{,}16nn-0{,}0154n^4&68660{,}7n-0{,}840n^3\\
53^\circ&-163539{,}8-467{,}79nn-0{,}0140n^4&67120{,}2n-0{,}936n^3\\
54^\circ&-274811{,}2-462{,}85nn-0{,}0127n^4&65559{,}1n-1{,}026n^3\\
55^\circ&-386100{,}9-457{,}34nn-0{,}0114n^4&63977{,}8n-1{,}109n^3.
\end{array}
\]
. . . . . .
\bigskip
\begin{center}\textsc{Gauss} to \textsc{Schumacher}.\quad Göttingen, 25 June 1831.\end{center}
. . . . . . As for calculating the coordinates, there are two problems to solve, namely:

1) Given the actual length \(R\) of a shortest line on the spheroid whose endpoints have coordinates \(x,y;\ x',y'\) respectively in the projection, find the distance [in] the projection, that is, the quantity \(r=\surd((x'-x)^2+(y'-y)^2)\). Here is the solution of the inverse problem. For all your cases the following is sufficiently accurate:
\[
\log R=\log r-k\,.\,\frac{(1-ee\sin\varphi^2)^2}{2aa(1-ee)}\,.\,\frac{yy+yy'+y'y'}3.
\]
Since \(y\) and \(y'\) need only be known approximately here, it is enough to make a preliminary calculation of the coordinates as though everything were in plano. The formula thus also serves to find \(r\) from \(R\). \(k\) is the modulus of Briggsian logarithms.

For the actual computation I put the formula in the following form:
\[
\log r=\log R+\{\alpha(y+y')^2+\beta(y-y')^2\}q,
\]
\noindent where
\[
\alpha=\frac{k}{4\,.\,206265},\quad\beta=\frac{k}{12\,.\,206265},\quad q=206265\,.\,\frac{(1-ee\sin\varphi^2)^2}{2aa(1-ee)};
\]
\par\noindent\rule{4cm}{.4pt}\par
{\footnotesize\noindent\hspace*{\parindent}[*) These formulas are found in Art. 13 at the bottom of p. 157.]\par}

% Printed216; principalPDF223.
\par\clearpage
\noindent216\hfill\textsc{Correspondence.}\hfill\mbox{}\par\medskip
\noindent here we have
\[
\begin{array}{l}
\log\alpha=3{,}72130-10\\
\log\beta=3{,}24418-10.
\end{array}
\]
For \(q\) I have an auxiliary table which need depend only on \(x\) and which I believe I have already sent you\textsuperscript{[*)]}. Clearly \(\alpha,\beta\) are independent of the unit of measure, whereas \(q\) is not. An auxiliary table for it depends both on the unit of measure and on the origin of the \(x\)-coordinates. If you use a different unit and a different origin, you would therefore have to transform my table first.

This gives you \(\log r-\log R\) as a decimal fraction. If you want the value directly in units of the seventh decimal place, you need only use \(\log\alpha=0{,}72130,\ \log\beta=0{,}24418\).

2) Given an angle measured at \(P\) between \(PP'\) and \(PP''\), find the angle corresponding to it in the projection. Here again the inverse problem takes its place.
\begin{center}
\begin{tikzpicture}[x=1cm,y=1cm,font=\normalsize]
\node at (1.8,2) {Spheroid};\node at (9,2) {Drawing in plano};
\draw (0,.7) circle(.035);\node[above left] at (0,.7) {$P''$};
\draw (2.1,1.2) circle(.035);\node[above right] at (2.1,1.2) {$P$};
\draw (2.85,.25) circle(.035);\node[above right] at (2.85,.25) {$P'$};
\draw (7.1,.7) circle(.035);\node[above left] at (7.1,.7) {$p''$};
\draw (9.3,1.2) circle(.035);\node[above right] at (9.3,1.2) {$p$};
\draw (10.05,.25) circle(.035);\node[above right] at (10.05,.25) {$p'$};
\end{tikzpicture}
\end{center}
Let the coordinates of \(p,p',p''\) be respectively
\[
\begin{array}{ccc}x,&x',&x'',\\y,&y',&y'',\end{array}
\]
\noindent then:
\[
\begin{aligned}
P'PP''={}&p'pp''-q^*(x'-x)\left(\frac{2y+y'}3\right)\\
&+q^{**}(x''-x)\left(\frac{2y+y''}3\right).
\end{aligned}
\quad\text{(gives the required result in arcseconds)}
\]
For \(q^*,q^{**}\) you can use the value of \(q\) that corresponds to the argument \(x\). If you want greater precision, the argument for \(q^*\) is \(\frac{2x+x'}3=x+\frac13(x'-x)\), and the argument for \(q^{**}\) is \(\frac{2x+x''}3=x+\frac13(x''-x)\). I always write \(\frac{2y+y'}3\) in the form \(y+\frac13(y'-y)\), etc. You see that the whole procedure then comes down to first adding to each measured azimuth (from
\par\noindent\rule{4cm}{.4pt}\par
{\footnotesize\noindent\hspace*{\parindent}[*) Page 210, where \(H=q\).]\par}

% Printed217; principalPDF224.
\par\clearpage
\noindent\mbox{}\hfill{\small\textsc{On the Formulas for the Hanoverian Land Survey.}}\hfill217\par\medskip
\noindent \(P\) to \(P'\)) the correction \(-q(x'-x)(y+\frac13(y'-y))\), which gives what I call the azimuth in plano. Here I use azimuth on the spheroid in a sense distinct from astronomical azimuth, which is the angle with the actual meridian. I mean the angle with a line parallel to the fundamental meridian, or, more strictly, with a line on the spheroid whose projection in plano is parallel to the axis of abscissas.

Thus the azimuth in plano is always \(\operatorname{arc\,tang}\frac{y'-y}{x'-x}\).

Once these corrections are applied, the calculation with all the angles proceeds as if everything were in plano. You need no instructions for calculating the coordinates from there. Clearly, an approximate knowledge of \(x,y,x'\), etc., also suffices for this computation, as obtained by calculating initially without corrections. If you want everything to agree to \(0''{,}001\), you can, if necessary, refine the calculation once \(x,y\), etc., are known more precisely.

You can readily see, moreover, that calculation 1) above, that is, the passage from \(R\) to \(r\), need only be carried out for one line (the baseline). Thereafter, knowledge of the individual lines on the spheroid is, generally speaking, of no interest unless you come to another baseline, where the reverse procedure (from \(r\) to \(R\)) is to be applied. In any event, in all your cases \(R\) and \(r\) are always very nearly equal, and the difference has no appreciable effect on centering the angles. In the westernmost part of Westphalia I did take it into account, but only for the honor of the calculation, since it makes no difference even there.
\begin{center}\textsc{Gauss} to \textsc{Schumacher}.\quad Göttingen, 9 December 1838.\end{center}
The requested formula is as follows\textsuperscript{[*)]}:
\[
\begin{aligned}
\varphi={}&\psi+(6ff+48f^4+426f^6+4080f^8\ldots)\sin2\psi\\
&\quad+(21f^4+336f^6+4264f^8\ldots)\sin4\psi\\
&\qquad+\left(\frac{302}3f^6+2416f^8\ldots\right)\sin6\psi\\
&\qquad\quad+\left(\frac{1097}2f^8\ldots\right)\sin8\psi\\
&\qquad\qquad\text{etc.}
\end{aligned}
\]
\par\noindent\rule{4cm}{.4pt}\par
{\footnotesize\noindent\hspace*{\parindent}[*) Cf. Art. 25, p. 179.]\par}
\noindent{\small IX.}\hfill28\hspace*{.12\linewidth}\par

% Printed218; principalPDF225.
\par\clearpage
\noindent218\hfill{\footnotesize\textsc{Correspondence.\quad On the Formulas for the Hanoverian Land Survey.}}\par\medskip
Since I have not worked on these matters for quite a few years, it took me a great deal of searching to find the formula again, and then more searching, or rather calculating, to determine the meaning of \(f\), which was not written on the sheet. It is \(f=\frac14e\), where \(e\) denotes the eccentricity of the generating ellipse. Or, if the flattening is \(\frac1n\), that is, \(\frac{\text{Earth's axis}}{\text{diam. of the equator}}=\frac{n-1}n\), then \(f=\frac{\surd(2n-1)}{4n}\). I used this \(f\) because using \(e\) would have made all the coefficients fractions, whereas here most of them are integers. I carried out the numerical calculation throughout with logarithms to 10 decimal places.
\begin{center}\rule{3cm}{.4pt}\end{center}
\bigskip
\begin{center}\textsc{Remarks.}\end{center}
\begingroup\footnotesize
The first two and the last of the preceding letters are printed from the originals held in the \textsc{Gauss} Archive. Copies made from the originals were used for the letters of 17 May and 25 June 1831. A few slips of the pen and minor errors in arithmetic have been corrected in the printed text. The letter of 9 December 1838 is a reply to \textsc{Schumacher}'s request for the derivation of the numerical formula for \(\varphi-\psi\) given in the letter of 18 April 1830. \textsc{Gauss}'s expression
\[
\varphi=\psi+520''{,}469\,3364\sin2\psi\ \text{etc.}
\]
\noindent contains an error in arithmetic, however, as the surviving sheet that \textsc{Gauss} used for the calculation also shows. It should read as given on p. 206:
\[
\varphi=\psi+520''{,}463\,3364\sin2\psi\ \text{etc.}
\]
\noindent This also makes the last digit in the table given there for \(\varphi-\psi\) uncertain. Furthermore, for Göttingen one obtains \(\psi^\circ=51^\circ23'20''{,}6082\), p. 207, in place of \(\ldots20''{,}6024\), as it reads in the original.
\par\noindent\hfill\textsc{Krüger, Börsch.}\par
\endgroup
\vspace{1.8cm}\begin{center}\rule{4cm}{.4pt}\end{center}

% Unnumbered title219; principalPDF226.
\par\clearpage\thispagestyle{empty}
\null\vspace*{.23\textheight}
\begin{center}{\Large TRIGONOMETRIC POINT DETERMINATION.}\end{center}
\vfill\noindent\hfill28\textsuperscript{*}\hspace*{.15\linewidth}\par\vspace*{.14\textheight}

% Blank220; principalPDF227.
\par\clearpage\null\thispagestyle{empty}

% Unnumbered opening221; principalPDF228.
\par\clearpage\thispagestyle{empty}
\vspace*{2.8\baselineskip}
\begin{center}{\Large\textls[160]{POSTHUMOUS PAPERS.}}\end{center}
\vspace{2\baselineskip}
\begin{center}[1.]\\[5pt]Final Result for the Position of a Point in a Plane Fixed\\by Intersection from Three Known Points.\end{center}
Let \(10,20,30\) denote the three observed directions [toward \(P\)] and \(\alpha,\beta,\gamma\) the corresponding distances.

Let the three individual results from the combinations \(2-3,\ 1-3,\ 1-2\) be \(A,B,C\), which are also the angles of the triangle formed by them, and let the sides opposite them be \(a,b,c\).
\begin{center}
\begin{tikzpicture}[x=1cm,y=1cm,line width=.6pt,font=\normalsize]
\coordinate (one) at (0,1.4);\coordinate (two) at (9.4,6.045);\coordinate (three) at (10.5,-.06);
\coordinate (B) at (4,2.8);\coordinate (C) at (7,3.85);\coordinate (A) at (5.25,2.25);\coordinate (P) at (5.44,2.86);
\draw (one)--(B)--(C);\draw (A)--(C)--(two);\draw (B)--(A)--(three);
\draw[line width=.4pt] (5.30155,3.25555)--(P)--(5.6573,2.62236) (P)--(5.1843,2.2789);
\foreach \p in {one,two,three,P}{\fill[white] (\p) circle(.045);\draw (\p) circle(.045);}
\node[above] at (one) {$1$};\node[right] at (two) {$2$};\node[right] at (three) {$3$};
\node[below] at (A) {$A$};\node[above left] at (B) {$B$};\node[right] at (C) {$C$};\node[above right] at (P) {$P$};
\node at (5.48,3.54) {$a$};\node at (6.46,2.92) {$b$};\node at (4.62,2.31) {$c$};
\node at (5.18,3.05) {$x$};\node at (5.49,2.61) {$y$};\node at (5.10,2.65) {$z$};
\end{tikzpicture}
\end{center}
